\section{Introduction}
\label{sec:introduction}

In an expander graph, the edge boundary of a set of vertices controls
its distance from the empty set or the whole vertex set. Coboundary
expansion extends this condition to sets of faces in a simplicial
complex. Over $\F$, the coboundary of a set of $i$-faces consists of
the $(i+1)$-faces containing an odd number of them. We construct
complexes of increasing size for which this boundary controls distance
from the space of coboundaries, uniformly in the size and with a bounded
number of top-dimensional faces at each vertex.

Let $Y$ be a finite pure $s$-dimensional simplicial complex, and let $Y(i)$
denote its faces with $i+1$ vertices. Purity means that every face is
contained in an $s$-face. Write $C^i(Y;\F)$ for the functions from
$Y(i)$ to $\F$. We use the weights
\begin{equation}\label{eq:weights}
 w_i^Y(\sigma)=
 \frac{\#\{\tau\in Y(s):\sigma\subseteq\tau\}}
 {\binom{s+1}{i+1}\,|Y(s)|},
 \qquad
 \norm{f}_{Y,i}=\sum_{\sigma:\,f(\sigma)\ne0}w_i^Y(\sigma).
\end{equation}
Equivalently, sample an $s$-face uniformly and then one of its $i$-faces
uniformly; $\norm{f}_{Y,i}$ is the probability that $f$ is nonzero there.
All simplicial chains and cochains have coefficients in $\F$ unless
specified otherwise. The coboundary is
\[
 (\delta_i f)(\tau)=\sum_{\substack{\sigma\subset\tau\\|\sigma|=i+1}}
 f(\sigma)\quad\text{in }\F.
\]
Write $Z^i(Y)=\ker\delta_i$ and
$B^i(Y)=\im\delta_{i-1}$, taking $B^0(Y)$ to be the constant functions.
For a subspace $A\subseteq C^i(Y;\F)$, put
\[
 \dist_{Y,i}(f,A)=\min_{a\in A}\norm{f-a}_{Y,i}.
\]
Thus $Z^0(Y)=B^0(Y)$ when $Y$ is connected.

\begin{theorem}\label{thm:main}
For every integer $d\ge3$, there exist constants $D<\infty$, $\eps>0$,
and finite connected pure $d$-dimensional simplicial complexes $X_m$,
with $|X_m(0)|\longrightarrow\infty$, such that every vertex belongs to
at most $D$ faces of dimension $d$ and, for every
$f\in C^i(X_m;\F)$,
\begin{equation}\label{eq:goal}
 \norm{\delta_i f}_{X_m,i+1}
 \ge \eps\,\dist_{X_m,i}(f,B^i(X_m))
 \qquad(0\le i<d).
\end{equation}
The constants are independent of $m$, $i$, and $f$.
\end{theorem}

The distinction between cocycles and coboundaries is essential. If
$f\in Z^i(Y)\setminus B^i(Y)$, its coboundary vanishes while its distance
from $B^i(Y)$ is positive. Thus \eqref{eq:goal} requires the vanishing of
$H^i(Y;\F)=Z^i(Y)/B^i(Y)$, with reduced cohomology understood in degree
zero. A uniform inequality against $Z^i(Y)$ alone allows nonzero
cohomology. Our proof supplies both that quantitative inequality and
the required vanishing for the same complexes.

\paragraph{Prior work.}
The higher-dimensional coisoperimetric viewpoint developed through the
random-complex work of Linial and Meshulam \cite{LM2006} and Meshulam and
Wallach \cite{MW2009}. Their models retain a complete lower skeleton and
study the disappearance of cohomology as top-dimensional faces are
added. Gromov connected cohomological filling inequalities with
topological overlap and asked for arbitrarily large complexes with
bounded vertex incidence and uniform filling bounds
\cite[Sections~2.3--2.5 and~2.14]{Gromov2010}. The full coboundary
expansion problem also demands vanishing cohomology; a recent formulation
is \cite[Problem~1.3]{CL2025}.

There are two different degree conditions in this history. Lubotzky and
Meshulam constructed random Latin-square coboundary expanders in
dimension two \cite{LM2015}; Lubotzky, Luria, and Rosenthal constructed
random Steiner-system coboundary expanders in every dimension
\cite{LLR2019}. These results bound the number of top-dimensional faces
through a codimension-one face. The Latin-square complexes have a
complete tripartite $1$-skeleton, and the Steiner-system complexes have
a complete $(d-1)$-skeleton. In dimensions at least two, vertex degrees
grow with the size. The degree bound in \cref{thm:main} is at vertices.

For bounded vertex degree, Kaufman, Kazhdan, and Lubotzky established
cosystolic and topological expansion for two-dimensional skeleta of
Ramanujan complexes \cite{KKL2016}. Evra and Kaufman obtained
bounded-degree cosystolic expanders in every dimension through a
local-to-global theorem \cite{EK2024}. The associated expansion methods
include the building-like complexes of Lubotzky, Meshulam, and Mozes
\cite{LMM2016}, the strongly symmetric complexes of Kaufman and Oppenheim
\cite{KO2021}, the coefficient-independent methods of Kaufman and Mass
\cite{KM2021}, and the quantitative local-to-global results of Dikstein
and Dinur \cite{DD2024}. We use only the distance-to-cocycles part of
cosystolic expansion, stated precisely in \eqref{eq:cosystolic}.

Chapman and Lubotzky proved the existence of bounded-degree
two-dimensional coboundary expanders over $\F$
\cite[Theorem~1.9]{CL2025}. In higher ambient dimensions, Kaufman,
Oppenheim, and Weinberger \cite{KOW2025}, and Oppenheim and
Valentiner-Branth \cite{OVBNonAbelian2025}, proved degree-one coboundary
expansion for coset complexes with suitable non-Abelian coefficients.
Those degree-one results and the cosystolic constructions above do not
give the vanishing in every degree required by \cref{thm:main}.
Together with the two-dimensional result and the graph case, our theorem
gives bounded-degree $\F$ coboundary expanders in every positive
dimension.

Our complexes belong to an established algebraic construction program.
Kaufman and Oppenheim developed locally spectral expanding coset
geometries \cite{KO2023}; O'Donnell and Pratt used Chevalley groups
\cite{OP2022}; and Grave de Peralta and Valentiner-Branth developed
Kac--Moody--Steinberg group constructions \cite{GdPVB2025}.
Oppenheim and Valentiner-Branth proved the cosystolic expansion theorem
for these KMS complexes that supplies our quantitative input \cite{OVB}.
Their affine-type-$A$ polynomial matrix construction
\cite[Example~4.13]{OVB} already uses the simple root groups appearing
here, with quotients by irreducible polynomials. Powers of $t$ also occur,
with different local groups, in the degree-one construction of
\cite[Section~2.5 and Theorem~6.3]{KOW2025}.
For the local groups used here, we pass to the tower of quotients modulo
$t^m$ and prove the cohomology vanishing needed to obtain coboundary
expansion in all degrees below $d$.

\paragraph{Proof strategy.}
Fix $d$ and put $n=2d$ and $N=2d+1$. Over a sufficiently large fixed
odd prime field $k$, let $U\le\SL_N(k[t])$ be the inverse image of the
upper unitriangular group under evaluation at zero. We construct an
infinite $n$-dimensional coset complex $L$ for $U$ and its finite
congruence quotients $T_m$. The desired $X_m$ is the $d$-skeleton of
$T_m$. The local groups are fixed finite groups, which gives uniform
degree bounds; the congruence quotients grow with $m$.

The first step is geometric. We identify $L$ with the subcomplex
opposite a chamber in an affine twin building. A filtration by
codistance expresses the relative chain groups of each attachment in
terms of a spherical opposition complex. The opposition-complex input
of \cite{OVB} makes these attachments acyclic in the required range.
Since the whole affine building is contractible, $L$ has vanishing
reduced homology below degree $n-1$. Its simplex stabilizers have odd
order, so its chain complex gives a partial projective resolution over
$\F[U]$. The orbit complex $U\backslash L$ is a single simplex, which
yields the corresponding homology vanishing for $U$.

The second step passes from $U$ to each congruence subgroup. Morrow's
pro Tor-unitality theorem \cite{Morrow} and Iwasa's homology pro-stability
theorem \cite{Iwasa} apply to the ideals $t^m k[t]$. Their initial
conclusion concerns an inverse system, whereas we need a statement at
every level. Odd-index transfer makes the homology maps towards
shallower levels surjective over $\F$. The cokernels of stabilization
are therefore both pro-zero and connected by surjective transition
maps, so each cokernel vanishes. Two stabilization steps then show that
ambient conjugation acts trivially in degrees $1,\ldots,d-1$.
Averaging over the finite odd-order quotient transfers the vanishing
from $U$ to the congruence subgroup, and hence to $T_m$.

Finally, the KMS theorem gives a uniform inequality against cocycles
on the $(n-1)$-skeleton of $T_m$. Its weighted norms agree with the
corresponding norms on $T_m$ because every panel lies in the same
number of chambers. Passing to the $d$-skeleton requires a further
comparison: the number of chambers through a $d$-face is bounded above
and below, and the normalization contributes the same factor in two
successive cochain degrees. This factor cancels in the expansion
inequality. The cohomology vanishing replaces cocycles by coboundaries
and completes the proof.

\paragraph{Organization.}
Section~\ref{sec:construction} defines the congruence complexes.
Section~\ref{sec:geometry} proves the local group identities and the
acyclicity of $L$. Section~\ref{sec:homology} proves the cohomology
vanishing of the finite quotients. Section~\ref{sec:expansion} checks
the KMS hypotheses, compares skeleton weights, and proves
\cref{thm:main}.
